\documentclass[10pt,a4paper]{amsart}
\usepackage{amsmath,amssymb,amsthm,graphicx}
% Retained preamble for the Stacks Project chapter {CH} (shared project preamble).
% The chapter file pinned-original-source.tex only carries \input{{preamble}} on its first line,
% so no preamble text lives inside the chapter itself.  The Stacks Project uses one single shared
% preamble for every chapter: the snapshot stages it three times (sources/stacks-chow/preamble.tex,
% sources/stacks-etale-cohomology/preamble.tex and focused-seeds/AGP/sources/stacks-snapshot-20261011/
% preamble.tex) and those copies are identical in size.  The block below is copied mechanically,
% line for line, from sources/stacks-chow/preamble.tex lines 14-18 and 154-217: the Xy-pic setup and
% the native theorem declarations together with the mathematical macros.  Nothing was added, removed
% or redefined; the wrapper class and the amsmath/amssymb/amsthm/graphicx base packages are supplied
% by the assembler.
% For commutative diagrams we use Xy-pic
% Chapter: stacks-cohomology.tex -- Cohomology of Sheaves.
\usepackage[all]{xy}

% We use 2cell for 2-commutative diagrams.
\xyoption{2cell}

% Theorem environments.
%
\theoremstyle{plain}
\newtheorem{theorem}[subsection]{Theorem}
\newtheorem{proposition}[subsection]{Proposition}
\newtheorem{lemma}[subsection]{Lemma}

\theoremstyle{definition}
\newtheorem{definition}[subsection]{Definition}
\newtheorem{example}[subsection]{Example}
\newtheorem{exercise}[subsection]{Exercise}
\newtheorem{situation}[subsection]{Situation}

\theoremstyle{remark}
\newtheorem{remark}[subsection]{Remark}
\newtheorem{remarks}[subsection]{Remarks}

\numberwithin{equation}{subsection}

% Macros
%
\def\lim{\mathop{\mathrm{lim}}\nolimits}
\def\colim{\mathop{\mathrm{colim}}\nolimits}
\def\Spec{\mathop{\mathrm{Spec}}}
\def\Hom{\mathop{\mathrm{Hom}}\nolimits}
\def\Ext{\mathop{\mathrm{Ext}}\nolimits}
\def\SheafHom{\mathop{\mathcal{H}\!\mathit{om}}\nolimits}
\def\SheafExt{\mathop{\mathcal{E}\!\mathit{xt}}\nolimits}
\def\Sch{\mathit{Sch}}
\def\Mor{\mathop{\mathrm{Mor}}\nolimits}
\def\Ob{\mathop{\mathrm{Ob}}\nolimits}
\def\Sh{\mathop{\mathit{Sh}}\nolimits}
\def\NL{\mathop{N\!L}\nolimits}
\def\CH{\mathop{\mathrm{CH}}\nolimits}
\def\proetale{{pro\text{-}\acute{e}tale}}
\def\etale{{\acute{e}tale}}
\def\QCoh{\mathit{QCoh}}
\def\Ker{\mathop{\mathrm{Ker}}}
\def\Im{\mathop{\mathrm{Im}}}
\def\Coker{\mathop{\mathrm{Coker}}}
\def\Coim{\mathop{\mathrm{Coim}}}

% Boxtimes
%
\DeclareMathSymbol{\boxtimes}{\mathbin}{AMSa}{"02}

%
% Macros for moduli stacks/spaces
%
\def\QCohstack{\mathcal{QC}\!\mathit{oh}}
\def\Cohstack{\mathcal{C}\!\mathit{oh}}
\def\Spacesstack{\mathcal{S}\!\mathit{paces}}
\def\Quotfunctor{\mathrm{Quot}}
\def\Hilbfunctor{\mathrm{Hilb}}
\def\Curvesstack{\mathcal{C}\!\mathit{urves}}
\def\Polarizedstack{\mathcal{P}\!\mathit{olarized}}
\def\Complexesstack{\mathcal{C}\!\mathit{omplexes}}
% \Pic is the operator that assigns to X its picard group, usage \Pic(X)
% \Picardstack_{X/B} denotes the Picard stack of X over B
% \Picardfunctor_{X/B} denotes the Picard functor of X over B
\def\Pic{\mathop{\mathrm{Pic}}\nolimits}
\def\Picardstack{\mathcal{P}\!\mathit{ic}}
\def\Picardfunctor{\mathrm{Pic}}
\def\Deformationcategory{\mathcal{D}\!\mathit{ef}}
\usepackage{cleveref}
\crefname{originallemma}{original lemma}{original lemmas}
\Crefname{originallemma}{Original lemma}{Original lemmas}
\makeatletter
\newlabel{derived-lemma-subcategory-right-acyclics}{{05T8}{source}{}{sourcepointer-1}{}}
\newlabel{derived-lemma-subcategory-right-acyclics@cref}{{[originallemma][0][]05T8}{[1][0][]source}{}{}{}}
\makeatother
\crefname{originallemma}{original lemma}{original lemmas}
\Crefname{originallemma}{Original lemma}{Original lemmas}
\makeatletter
\newlabel{injectives-lemma-abelian-sheaves-space}{{01DG}{source}{}{sourcepointer-2}{}}
\newlabel{injectives-lemma-abelian-sheaves-space@cref}{{[originallemma][0][]01DG}{[1][0][]source}{}{}{}}
\makeatother
\title{Flasque modules are acyclic for the global section functor and for sections over any open}
\author{The Stacks Project Authors}
\date{stacks-project:cohomology.tex:local-snapshot-20261011; GNU FDL 1.2 or later}
\begin{document}
\maketitle
\noindent\textit{Editorial note.} All delivered mathematical text is the author's own, copied from the pinned Stacks Project chapter file cohomology.tex; no statement or proof step was rewritten, reordered or removed.\par\medskip
\noindent\textit{Editorial note.} The proof is complete as written: the maximality argument, the correction using flasqueness of the subsheaf and the final gluing are all present in the source.\par\medskip
\noindent\textit{Editorial note.} The counted spans carry 2 navigation stubs and no bibliography citation.\par\medskip
\noindent\textit{Editorial note.} The preamble is the chapter's retained minimal preamble (sources/stacks-cohomology/retained-preamble.tex lines 1-81), a verbatim excerpt of the native preamble.\par\medskip
\noindent\textit{Editorial note.} The label lemma-flasque-acyclic is defined only in cohomology.tex, so no cross-chapter restatement had to be weighed for this card.\par\medskip
\noindent\textit{Allowed standard prerequisites (editorial).}\begin{itemize}
\item Criterion for a class of objects to be right acyclic. Derived Categories, Lemma derived-lemma-subcategory-right-acyclics, invoked at proof lines 1920-1921. Supplies that the two verified conditions on the class of flasque modules imply acyclicity.
\item Every O\_X-module embeds into a flasque module. Injectives, Lemma injectives-lemma-abelian-sheaves-space, invoked at proof lines 1922-1924. Supplies the embedding condition, since injective modules are flasque.
\item The sheaf condition on G for a union of compatible sections. Standard sheaf theory, used at proof lines 1941-1943 and 1953-1955. Supplies the existence of the glued section over a union of opens.
\end{itemize}\medskip
\ifdefined\hypertarget\hypertarget{sourcepointer-1}{}\fi
\noindent\textit{Source reference.} Original lemma 05T8: sources/stacks-chow/tags.txt, source lines 7472--7472. Pointer outside the counted span, invoked once for context; the load-bearing argument of the counted statement is the author's proof reproduced in full in the proof span.\par\medskip
\ifdefined\hypertarget\hypertarget{sourcepointer-2}{}\fi
\noindent\textit{Source reference.} Original lemma 01DG: sources/stacks-chow/tags.txt, source lines 1797--1797. Pointer outside the counted span, invoked once for context; the load-bearing argument of the counted statement is the author's proof reproduced in full in the proof span.\par\medskip
\begin{lemma}
\label{lemma-flasque-acyclic}
Let $(X, \mathcal{O}_X)$ be a ringed space. Any flasque $\mathcal{O}_X$-module
is acyclic for $R\Gamma(X, -)$ as well as $R\Gamma(U, -)$ for any
open $U$ of $X$.
\end{lemma}

\begin{proof}
We will prove this using
Derived Categories, Lemma \ref{derived-lemma-subcategory-right-acyclics}.
Since every injective module is flasque we see that we can embed
every $\mathcal{O}_X$-module into a flasque module, see
Injectives, Lemma \ref{injectives-lemma-abelian-sheaves-space}.
Thus it suffices to show that given a short exact sequence
$$
0 \to \mathcal{F} \to \mathcal{G} \to \mathcal{H} \to 0
$$
with $\mathcal{F}$, $\mathcal{G}$ flasque, then $\mathcal{H}$
is flasque and the sequence remains short exact after taking sections
on any open of $X$. In fact, the second statement implies the first.
Thus, let $U \subset X$ be an open subspace. Let $s \in \mathcal{H}(U)$.
We will show that we can lift $s$ to a section of $\mathcal{G}$
over $U$. To do this consider the set $T$ of pairs $(V, t)$
where $V \subset U$ is open and $t \in \mathcal{G}(V)$ is a section
mapping to $s|_V$ in $\mathcal{H}$.
We put a partial ordering on $T$ by setting
$(V, t) \leq (V', t')$ if and only if $V \subset V'$ and $t'|_V = t$.
If $(V_\alpha, t_\alpha)$, $\alpha \in A$
is a totally ordered subset of $T$, then $V = \bigcup V_\alpha$
is open and there is a unique section $t \in \mathcal{G}(V)$
restricting to $t_\alpha$ over $V_\alpha$ by the sheaf condition on
$\mathcal{G}$. Thus by Zorn's lemma there exists a maximal element
$(V, t)$ in $T$. We will show that $V = U$ thereby finishing the proof.
Namely, pick any $x \in U$. We can find a small open neighbourhood
$W \subset U$ of $x$ and $t' \in \mathcal{G}(W)$ mapping to $s|_W$
in $\mathcal{H}$. Then $t'|_{W \cap V} - t|_{W \cap V}$ maps to
zero in $\mathcal{H}$, hence comes from some section
$r' \in \mathcal{F}(W \cap V)$. Using that $\mathcal{F}$ is flasque
we find a section $r \in \mathcal{F}(W)$ restricting to $r'$
over $W \cap V$. Modifying $t'$ by the image of $r$ we may
assume that $t$ and $t'$ restrict to the same section over
$W \cap V$. By the sheaf condition of $\mathcal{G}$
we can find a section $\tilde t$ of $\mathcal{G}$ over
$W \cup V$ restricting to $t$ and $t'$.
By maximality of $(V, t)$ we see that $V \cup W = V$.
Thus $x \in V$ and we are done.
\end{proof}

\end{document}
